%% ═══════════════════════════════════════════════════════════════════════════
\section{Lecture 6: Diffusion, Heat Conduction, Viscosity, and the Reynolds Number}
\label{sec:lec06}
%% ═══════════════════════════════════════════════════════════════════════════
\date{14/05/2026}

This lecture unifies three apparently different physical phenomena — particle
diffusion, heat conduction, and viscous flow — by showing that they all arise
from the \emph{same microscopic mechanism}: particles random-walking and
transporting whatever scalar or vector quantity they carry. Dimensional analysis
then fixes the form of each transport coefficient, and dimensionless ratios of
these coefficients encode deep physics about the dominant carrier.

%%───────────────────────────────────────────────────────────────────────────
\subsection{What Is Diffusion?}
\label{sec:lec06:diffusion-intro}
%%───────────────────────────────────────────────────────────────────────────

\textbf{The Physical Picture:}
Diffusion is the net transport of a quantity (particle number, heat, momentum)
driven by a gradient, through the statistical bias of random walkers. If there
are more ink molecules on one side of an imaginary surface than on the other,
more will cross from the dense side per unit time than from the dilute side,
producing a net flux even though individual particles move randomly. The same
picture applies to heat (particles carry thermal energy) and to momentum
(particles carry bulk velocity).

Canonical examples:
\begin{itemize}
    \item Dye in a liquid: ink diffuses until concentration is uniform.
    \item Heat: Fourier's law—heat flows from hot to cold.
    \item Momentum: viscosity—fast fluid layers drag slow ones.
    \item Social/biological systems: epidemic spreading, nutrient transport.
\end{itemize}

%%───────────────────────────────────────────────────────────────────────────
\subsection{Fick's Law from Dimensional Analysis}
\label{sec:lec06:fick-derivation}
%%───────────────────────────────────────────────────────────────────────────

\subsubsection*{Physical Motivation}
We want to find the particle flux $\mathbf{J}$ (number of particles crossing
unit area per unit time) in terms of the fundamental microscopic quantities.
The relevant parameters are:
\begin{itemize}
    \item $n$ — local particle number density, $[n] = \si{m^{-3}}$,
    \item $\ell$ — mean free path (typical distance a particle travels
          before a collision resets its direction), $[\ell] = \si{m}$,
    \item $v$ — typical thermal speed, $[v] = \si{m\,s^{-1}}$.
\end{itemize}
The quantity to find has dimensions $[J] = \si{m^{-2}\,s^{-1}}$.

\subsubsection*{Dimensional counting}
Build $J$ from $n$, $\ell$, and $v$. There are three variables and three
independent base dimensions (m, s, and the pure number count embedded in $n$),
so Buckingham's theorem leaves one dimensionless combination,
\begin{equation}
    \Pi = n\,\ell^3 ,
    \label{eq:lec06:pi-group-diffusion}
\end{equation}
and the most general dimensional relation is
\begin{equation}
    J = n\,v\,\ell\;\Phi\!\left(n\ell^3\right) ,
    \label{eq:lec06:J-general}
\end{equation}
where $\Phi$ is an unknown function. The factor $nv\ell$ has units
$\si{m^{-3}}\cdot\si{m\,s^{-1}}\cdot\si{m} = \si{m^{-2}s^{-1}}$ as required.

\subsubsection*{Linearity selects the leading term}
When the number density is small (dilute limit, $n\ell^3 \ll 1$), particles
interact rarely and the flux must be linear in $n$: doubling the number of
particles doubles the flux. This forces $\Phi = \text{const}$, giving
\begin{equation}
    J \sim n\,v\,\ell .
    \label{eq:lec06:J-linear}
\end{equation}
More carefully, the flux is the imbalance between particles crossing a surface
from left and right. The particles that reach the surface from the left started
roughly one mean free path $\ell$ away:
\begin{equation}
    J_+ - J_- \approx n\!\left(x-\tfrac{\ell}{2}\right)\,\tfrac{v}{2}
    - n\!\left(x+\tfrac{\ell}{2}\right)\,\tfrac{v}{2}
    = -\frac{v\ell}{2}\,\frac{\partial n}{\partial x} .
    \label{eq:lec06:flux-imbalance}
\end{equation}
Collecting factors of order unity gives Fick's first law:
\begin{equation}
    \boxed{
    J = -D\,\nabla n,
    \qquad D \sim \ell\,v .
    }
    \label{eq:lec06:ficks-first-law}
\end{equation}
\textit{Physical significance:} The diffusion coefficient $D$ has units
$\si{m^2\,s^{-1}}$; it sets the rate at which a concentration pattern spreads
in space. The minus sign reflects the physical fact that transport is down the
gradient.

%%───────────────────────────────────────────────────────────────────────────
\subsection{Estimating the Diffusion Coefficient in Air}
\label{sec:lec06:D-estimate}
%%───────────────────────────────────────────────────────────────────────────

\subsubsection*{Physical Motivation}
To make $D \sim \ell v$ concrete, we need the mean free path $\ell$ and
thermal speed $v$ of air molecules at room temperature.

\subsubsection*{Thermal speed}
The Maxwell-Boltzmann distribution for one component of velocity gives
\begin{equation}
    v_{\rm th} \sim \sqrt{\frac{k_B T}{m}} .
    \label{eq:lec06:thermal-speed}
\end{equation}
For air (average molecular mass $\bar m \approx 29\,\si{u} \approx
4.8\times10^{-26}\,\si{kg}$) at $T = 300\,\si{K}$:
\begin{equation}
    v_{\rm th} \sim \sqrt{\frac{1.38\times10^{-23}\times300}{4.8\times10^{-26}}}
    \simeq 300\text{--}500\,\si{m\,s^{-1}} .
    \label{eq:lec06:vthermal-air}
\end{equation}
In practice one takes $v_{\rm th}\simeq 340\text{--}500\,\si{m\,s^{-1}}$
(speed of sound is the lower bound for a related quantity; the true rms speed
is somewhat higher).

\subsubsection*{Mean free path}
A particle sweeps out a cylinder of cross-sectional area $\sigma$ (the
collision cross section). It collides when another particle lies inside this
cylinder. The mean free path is therefore
\begin{equation}
    \ell = \frac{1}{\sqrt{2}\,n\,\sigma} ,
    \label{eq:lec06:mean-free-path}
\end{equation}
where the $\sqrt{2}$ accounts for relative velocities in a moving gas.
Estimating the cross section from the molecular size $d \sim 3\,\si{\AA} =
3\times10^{-10}\,\si{m}$:
\begin{equation}
    \sigma \sim \pi\!\left(\frac{d}{2}\right)^{\!2}
    \approx \pi\times(1.5\times10^{-10})^2
    \approx 7\times10^{-20}\,\si{m^2} .
    \label{eq:lec06:sigma-molecule}
\end{equation}
The number density of air at sea level follows from $n = P/(k_B T)$ or from
the ratio to water:
\begin{equation}
    n_{\rm air}
    = \frac{P_0}{k_B T}
    \approx \frac{10^5\,\si{Pa}}{1.38\times10^{-23}\,\si{J\,K^{-1}}\times300\,\si{K}}
    \approx 2.4\times10^{25}\,\si{m^{-3}} .
    \label{eq:lec06:nair}
\end{equation}
Alternatively, water has $n_{\rm water} = \rho_w N_A / M_{\rm mol} \approx
3.3\times10^{28}\,\si{m^{-3}}$, and since $\rho_{\rm air}/\rho_w \approx
1.3\times10^{-3}$, we get $n_{\rm air} \approx 4\times10^{25}\,\si{m^{-3}}$
— consistent to within a factor of two. Therefore
\begin{equation}
    \ell
    \approx \frac{1}{\sqrt{2}\times2.4\times10^{25}\times7\times10^{-20}}
    \approx \frac{1}{\sqrt{2}\times1.7\times10^6\,\si{m^{-1}}}
    \approx 4\times10^{-7}\,\si{m} \approx 100\,\si{nm} .
    \label{eq:lec06:mfp-air}
\end{equation}
(The real value is $\approx 70\,\si{nm}$ at $20^\circ\si{C}$, consistent with
our estimate.) The diffusion coefficient is therefore
\begin{equation}
    \boxed{
    D_{\rm air} \sim \ell\,v_{\rm th}
    \approx 10^{-7}\,\si{m}\times5\times10^{2}\,\si{m\,s^{-1}}
    \approx 5\times10^{-5}\,\si{m^2\,s^{-1}} .
    }
    \label{eq:lec06:Dair}
\end{equation}
\textit{Physical significance:} The real value for air is
$D \approx 2\times10^{-5}\,\si{m^2\,s^{-1}}$, in good order-of-magnitude
agreement. The estimate succeeded using only molecular size and temperature.

%%───────────────────────────────────────────────────────────────────────────
\subsection{The Diffusion Equation and Characteristic Timescales}
\label{sec:lec06:diffusion-equation}
%%───────────────────────────────────────────────────────────────────────────

\subsubsection*{Physical Motivation}
If particles flow according to Fick's first law, the continuity equation
$\partial n/\partial t + \nabla\cdot \mathbf{J} = 0$ immediately gives the
diffusion equation (Fick's second law):
\begin{equation}
    \boxed{
    \frac{\partial n}{\partial t} = D\,\nabla^2 n .
    }
    \label{eq:lec06:diffusion-equation}
\end{equation}
\textit{Physical significance:} The sign of $\nabla^2 n$ is negative at a
concentration peak (divergence of flux is positive there — more leaves than
arrives) so the equation says peaks erode and troughs fill, driving the system
toward uniformity.

\subsubsection*{Characteristic timescale by dimensional analysis}
The diffusion equation contains only two scales: the diffusion coefficient $D$
($[\si{m^2\,s^{-1}}]$) and the length scale $L$ of the initial pattern. The
only combination with units of time is
\begin{equation}
    \boxed{
    \tau_{\rm diff} \sim \frac{L^2}{D} .
    }
    \label{eq:lec06:diffusion-timescale}
\end{equation}
\textit{Physical significance:} The quadratic dependence on $L$ is
characteristic of diffusion: doubling the length scale quadruples the spread
time. This makes diffusion extremely slow over macroscopic distances.

\subsubsection*{Example: Perfume in a room}
How long does it take for perfume molecules to diffuse across a room of size
$L \sim 1\,\si{m}$?
\begin{equation}
    \tau_{\rm diff}
    \sim \frac{(1\,\si{m})^2}{2\times10^{-5}\,\si{m^2\,s^{-1}}}
    \approx 5\times10^4\,\si{s}
    \approx 14\,\si{hours} .
    \label{eq:lec06:perfume-molecular}
\end{equation}
Yet in practice you smell perfume within a second or two. The discrepancy is
large: a factor of $\sim 10^4$.

\subsubsection*{Turbulent diffusion}
The resolution is that everyday airflow is \emph{turbulent}. Macroscopic
eddies — stirring motions with effective size $L_{\rm turb}$ and turnover
speed $v_{\rm turb}$ — mix the gas on timescale $\tau_{\rm turb} \sim
L_{\rm turb}/v_{\rm turb}$. Defining an effective turbulent diffusivity:
\begin{equation}
    D_{\rm turb} \sim v_{\rm turb}\,L_{\rm turb} .
    \label{eq:lec06:Dturb}
\end{equation}
To smell perfume in $\tau \sim 1\,\si{s}$ across $L \sim 1\,\si{m}$, we need
\begin{equation}
    \boxed{
    D_{\rm turb} \sim \frac{L^2}{\tau} \sim \frac{1\,\si{m^2}}{1\,\si{s}}
    \approx 1\,\si{m^2\,s^{-1}} ,
    }
    \label{eq:lec06:Dturb-estimate}
\end{equation}
which is five orders of magnitude larger than the molecular $D_{\rm air}$.
\textit{Physical significance:} Turbulent mixing completely dominates
molecular diffusion at human scales in air. In the atmosphere, diffusion is
also anisotropic: horizontal eddy diffusion is far more efficient than
vertical diffusion because the vertical density gradient strongly resists
upward displacement.

%%───────────────────────────────────────────────────────────────────────────
\subsection{Heat Conduction from the Same Mechanism}
\label{sec:lec06:heat-conduction}
%%───────────────────────────────────────────────────────────────────────────

\subsubsection*{Physical Motivation}
Each gas molecule carries kinetic energy $\varepsilon \sim k_B T$ (per degree
of freedom). When a molecule random-walks from a hot region into a cold one,
it deposits energy. The heat flux is therefore the \emph{energy content per
particle} times the particle flux.

\subsubsection*{Derivation of Fourier's law}
Let $\varepsilon(T)$ be the internal energy per particle. The heat flux is
\begin{align}
    \mathbf{J}_Q
    &= \varepsilon(T)\,\mathbf{J}
    = -D\,\varepsilon(T)\,\nabla n .
    \label{eq:lec06:Jq-step1}
\end{align}
If the number density is uniform but the temperature varies,
$\nabla n \approx 0$ but $\nabla(n\varepsilon) = n\,\nabla\varepsilon
= n\,\frac{d\varepsilon}{dT}\,\nabla T$. Therefore
\begin{equation}
    \mathbf{J}_Q
    = -D\,n\,\frac{d\varepsilon}{dT}\,\nabla T
    = -D\,c_V\,\nabla T ,
    \label{eq:lec06:Jq-step2}
\end{equation}
where $c_V = n\,d\varepsilon/dT$ is the heat capacity per unit volume. This
is Fourier's law with thermal conductivity $\kappa = D\,c_V$:
\begin{equation}
    \boxed{
    \mathbf{J}_Q = -\kappa\,\nabla T,
    \qquad
    \kappa = D\,c_V .
    }
    \label{eq:lec06:fourier}
\end{equation}
\textit{Physical significance:} The thermal conductivity is not an independent
property — it is the diffusion coefficient times the heat capacity per volume.
For ideal gases, this means $\kappa \sim \ell\,v\,n\,k_B$.

\subsubsection*{The heat diffusion equation}
Applying continuity to the energy density $u = c_V T$:
\begin{equation}
    \rho\,c_p\,\frac{\partial T}{\partial t}
    = -\nabla\cdot\mathbf{J}_Q
    = \kappa\,\nabla^2 T ,
    \label{eq:lec06:heat-eq-derivation}
\end{equation}
which gives
\begin{equation}
    \boxed{
    \frac{\partial T}{\partial t}
    = D_{\rm thermal}\,\nabla^2 T,
    \qquad
    D_{\rm thermal} = \frac{\kappa}{\rho\,c_p} .
    }
    \label{eq:lec06:heat-equation}
\end{equation}
\textit{Physical significance:} The thermal diffusivity $D_{\rm thermal}$ has
the same units and the same physical origin as the particle diffusion
coefficient. For an ideal monatomic gas, $c_V \sim \frac{3}{2}nk_B$ and
$\rho c_p \sim \frac{5}{2}nk_B$, so $D_{\rm thermal} = \frac{3}{5}D
\sim D$. The small numeric factor confirms that heat spreads on the same
timescale as particle density for ideal gases.

Note that $c_p$ appeared instead of $c_V$ because work done on adjacent fluid
elements must be included (constant-pressure process). The ratio $\gamma =
c_p/c_V$ appears, but it is $O(1)$ and does not change the order-of-magnitude
conclusion.

%%───────────────────────────────────────────────────────────────────────────
\subsection{Momentum Diffusion and Viscosity}
\label{sec:lec06:viscosity}
%%───────────────────────────────────────────────────────────────────────────

\subsubsection*{Physical Motivation}
Gas molecules also carry bulk momentum. In a shear flow where the $x$-velocity
$v_x$ depends on $y$, a molecule crossing from a fast layer at $y-\ell$ to a
slow layer at $y$ carries extra $x$-momentum and deposits it, tending to
accelerate the slow layer. This is viscosity — friction from momentum diffusion.

\subsubsection*{Viscous stress and kinematic viscosity}
The momentum flux (stress) in the $x$-direction through a surface of constant
$y$ is, by the same mean-free-path argument:
\begin{equation}
    \sigma_{xy}
    = -\nu\,\rho\,\frac{\partial v_x}{\partial y},
    \qquad
    \nu \sim \ell\,v_{\rm th},
    \label{eq:lec06:viscous-stress}
\end{equation}
where $\nu = \eta/\rho$ is the kinematic viscosity ($[\nu] = \si{m^2\,s^{-1}}$).
The dynamic viscosity is $\eta = \rho\,\nu$.
\begin{equation}
    \boxed{
    \nu \sim \ell\,v_{\rm th} \sim D_{\rm particle} .
    }
    \label{eq:lec06:nu-estimate}
\end{equation}
\textit{Physical significance:} For an ideal gas, kinematic viscosity,
particle diffusivity, and thermal diffusivity are all set by the \emph{same}
product $\ell v_{\rm th}$. They differ only by numerical prefactors of
order unity.

\subsubsection*{The Navier-Stokes equation}
Applying momentum conservation to a fluid element and adding the viscous
diffusion term gives the Navier-Stokes equation:
\begin{equation}
    \boxed{
    \rho\left(\frac{\partial \mathbf{v}}{\partial t}
    + (\mathbf{v}\cdot\nabla)\mathbf{v}\right)
    = -\nabla P + \eta\,\nabla^2\mathbf{v} + \rho\,\mathbf{g} .
    }
    \label{eq:lec06:navier-stokes}
\end{equation}
\begin{itemize}
    \item The $(\mathbf{v}\cdot\nabla)\mathbf{v}$ term is the \emph{advective}
    (convective) derivative — it is nonlinear in $\mathbf{v}$ and
    responsible for turbulence.
    \item The $\eta\nabla^2\mathbf{v}$ term is the viscous diffusion of
    momentum.
    \item The pressure gradient $-\nabla P$ drives flow from high to low
    pressure.
\end{itemize}
\textit{Physical significance:} The richness of fluid dynamics comes entirely
from the nonlinear advection term, which the viscous term opposes. Their ratio
is the Reynolds number.

%%───────────────────────────────────────────────────────────────────────────
\subsection{The Prandtl Number}
\label{sec:lec06:prandtl}
%%───────────────────────────────────────────────────────────────────────────

\subsubsection*{Physical Motivation}
We have three transport coefficients, all with the same units: $D$ (particle
diffusivity), $D_{\rm thermal}$ (thermal diffusivity), and $\nu$ (kinematic
viscosity). For an ideal gas they are all $\sim \ell v_{\rm th}$. But in a
liquid metal or a polymer melt the dominant carriers of heat and momentum are
completely different. The Prandtl number encodes this.

\begin{equation}
    \boxed{
    \mathrm{Pr} = \frac{\nu}{D_{\rm thermal}}
    = \frac{\text{kinematic viscosity}}{\text{thermal diffusivity}} .
    }
    \label{eq:lec06:prandtl}
\end{equation}
\textit{Physical significance:} $\mathrm{Pr}$ compares how efficiently
momentum versus heat is diffused.

\subsubsection*{Air ($\mathrm{Pr} \sim 1$)}
Both momentum and heat are carried by the same air molecules via the same
mean-free-path mechanism. Therefore $\nu \sim D_{\rm thermal}$ and
$\mathrm{Pr}_{\rm air} \approx 0.71$ — of order unity, as expected.

\subsubsection*{Mercury ($\mathrm{Pr} \ll 1$)}
Mercury is a liquid metal. Its atoms carry momentum through inter-atomic
collisions (same mechanism as gases, but in a dense liquid). However, heat is
mainly carried by \emph{conduction electrons}, which are far lighter, faster,
and more numerous than the atoms. Electrons diffuse heat much more efficiently
than the heavy atoms diffuse momentum:
\begin{equation}
    D_{\rm thermal}^{\rm Hg} \gg \nu^{\rm Hg}
    \qquad\Rightarrow\qquad
    \mathrm{Pr}_{\rm Hg} \approx 0.025 \ll 1 .
    \label{eq:lec06:pr-mercury}
\end{equation}
The general rule: \emph{anything that conducts electricity well (mobile
electrons) will also conduct heat well}, because the electron mean free path
is long and their speed is the Fermi velocity ($\sim 10^6\,\si{m\,s^{-1}}$),
not the thermal velocity. This is the Wiedemann-Franz law.

\subsubsection*{Long-chain polymers ($\mathrm{Pr} \gg 1$)}
A tangle of polymer chains transmits momentum very efficiently: pulling one
end of a chain exerts a force at the far end along the chain backbone, giving
extremely high viscosity. Heat transport, by contrast, still requires molecular
collisions and is slow. Consequently
\begin{equation}
    \nu^{\rm polymer} \gg D_{\rm thermal}^{\rm polymer}
    \qquad\Rightarrow\qquad
    \mathrm{Pr}_{\rm polymer} \gg 1 .
    \label{eq:lec06:pr-polymer}
\end{equation}

\begin{center}
\begin{tabular}{lcc}
    \toprule
    Material & $\mathrm{Pr}$ & Dominant heat carrier \\
    \midrule
    Air (ideal gas) & $\approx 0.7$ & Molecules \\
    Water & $\approx 7$ & Molecules (H-bond network) \\
    Mercury & $\approx 0.025$ & Conduction electrons \\
    Polymer melt & $\gg 10^2$ & No efficient carrier \\
    \bottomrule
\end{tabular}
\end{center}

%%───────────────────────────────────────────────────────────────────────────
\subsection{Drag on a Sphere and the Physical Meaning of Reynolds Number}
\label{sec:lec06:reynolds}
%%───────────────────────────────────────────────────────────────────────────

\subsubsection*{Physical Motivation}
A sphere of radius $R$ moves at speed $V$ through a fluid of density $\rho$
and kinematic viscosity $\nu$. We want the drag force $F$ using dimensional
analysis, and then connect $\mathrm{Re}$ to the competition between two
timescales.

\subsubsection*{Dimensional analysis of drag}
The variables are $F$ (force, $[MLT^{-2}]$), $\rho$ ($[ML^{-3}]$), $V$
($[LT^{-1}]$), $R$ ($[L]$), and $\nu$ ($[L^2T^{-1}]$). That is $N=5$
variables and $M=3$ independent base dimensions, so there are two
dimensionless groups. Choose
\begin{equation}
    \Pi_1 = \frac{F}{\rho V^2 R^2},
    \qquad
    \Pi_2 = \mathrm{Re} = \frac{V R}{\nu} .
    \label{eq:lec06:drag-pi-groups}
\end{equation}
Buckingham's theorem then gives
\begin{equation}
    \boxed{
    F = \rho V^2 R^2\,f(\mathrm{Re}) ,
    }
    \label{eq:lec06:drag-general}
\end{equation}
where $f$ is an unknown function that experiment or theory must supply.
\textit{Physical significance:} The drag coefficient $C_D \equiv F/(\rho V^2
A)$ is a function of Re alone (for a given body shape). All drag data
collapse onto a single curve when plotted as $C_D$ vs.\ Re.

\subsubsection*{Stokes regime ($\mathrm{Re} \ll 1$)}
At low Reynolds number, the fluid has time to adjust smoothly and viscosity
dominates. The problem is linear in $V$, so the force must be linear in $V$.
This is only consistent with Eq.~\eqref{eq:lec06:drag-general} if
\begin{equation}
    f(\mathrm{Re}) \propto \frac{1}{\mathrm{Re}}
    \quad\Rightarrow\quad
    F \propto \rho V^2 R^2 \cdot \frac{\nu}{VR}
    = \rho\nu VR = \eta VR .
    \label{eq:lec06:stokes-scaling}
\end{equation}
The exact solution (Stokes, 1851) gives the numerical prefactor:
\begin{equation}
    \boxed{
    F_{\rm Stokes} = 6\pi\,\eta\,V\,R .
    }
    \label{eq:lec06:stokes-drag}
\end{equation}
\textit{Physical significance:} Stokes drag is linear in speed and viscosity,
and linear in size — very different from turbulent drag ($F\propto V^2 R^2$).
Small organisms (bacteria, pollen) and falling droplets live in the Stokes
regime.

\subsubsection*{Physical interpretation of the Reynolds number}
Two natural timescales appear in the problem:
\begin{align}
    \tau_{\rm cross}
    &= \frac{R}{V}
    \quad\text{(time for fluid to flow past the object)},
    \label{eq:lec06:tau-cross}\\
    \tau_{\rm diff}
    &= \frac{R^2}{\nu}
    \quad\text{(time for momentum to diffuse a distance } R\text{)}.
    \label{eq:lec06:tau-diff}
\end{align}
Their ratio is exactly the Reynolds number:
\begin{equation}
    \boxed{
    \mathrm{Re}
    = \frac{\tau_{\rm diff}}{\tau_{\rm cross}}
    = \frac{R^2/\nu}{R/V}
    = \frac{VR}{\nu} .
    }
    \label{eq:lec06:re-timescale}
\end{equation}
\textit{Physical significance:}
\begin{itemize}
    \item \textbf{Re $\ll 1$}: viscous diffusion of momentum acts quickly
    compared with advection. The fluid ``knows'' the body is there before
    the fluid element reaches it — smooth, reversible Stokes flow.
    \item \textbf{Re $\gg 1$}: the fluid element sweeps past the body before
    momentum diffusion can communicate the body's presence to nearby
    fluid. Advection dominates; flow separates and turbulence develops.
\end{itemize}

%%───────────────────────────────────────────────────────────────────────────
\subsection{Summary of the Unified Transport Picture}
\label{sec:lec06:summary}
%%───────────────────────────────────────────────────────────────────────────

\textbf{The Physical Picture:}
All three classical transport phenomena — diffusion, heat conduction, and
viscosity — share a common mechanism: microscopic carriers (molecules,
electrons, phonons) performing random walks and transporting conserved
quantities across gradients. Dimensional analysis immediately gives the
transport coefficients; the Prandtl number then reveals whether one mechanism
dominates or whether different carriers are responsible for different
quantities.

\begin{itemize}
    \item Particle diffusion: $J = -D\nabla n$, \quad $D \sim \ell\,v_{\rm th}$,
    \quad $\partial_t n = D\nabla^2 n$.
    \item Heat conduction: $J_Q = -\kappa\nabla T$, \quad
    $\kappa = D\,c_V$, \quad $\partial_t T = D_{\rm th}\nabla^2 T$.
    \item Viscous flow: $\sigma_{xy} = -\eta\,\partial_y v_x$, \quad
    $\nu = \eta/\rho \sim \ell\,v_{\rm th}$, \quad full Navier-Stokes
    equation (Eq.~\eqref{eq:lec06:navier-stokes}).
    \item Diffusion timescale: $\tau \sim L^2/D$ (quadratic in $L$ — very slow
    at large scales, replaced by turbulence in real systems).
    \item Prandtl number: $\mathrm{Pr} = \nu/D_{\rm th}$; air $\approx 0.7$,
    mercury $\approx 0.025$, polymers $\gg 1$.
    \item Reynolds number: $\mathrm{Re} = VR/\nu = \tau_{\rm diff}/\tau_{\rm
    cross}$; low Re gives Stokes flow $F = 6\pi\eta VR$, high Re gives
    turbulent drag $F \sim \rho V^2 R^2$.
\end{itemize}
